\section{Limitations and conclusion}
\label{sec:conclusion}

\paragraph{Limitations.} All comparisons are single runs, and several change one setting on top of recipes that differ from each other. The downstream backbone predates the fixed-GSD stems, the dense context loss and the convolutional stems, so the current recipe is not yet evaluated downstream \TODO{final checkpoint}. The cross-sample gap cannot separate content from location or product identity, and it compares only runs whose targets are alike. The mechanism of \cref{eq:decomp} is supported by probes and ablations but not proven. Polar regions and the 1\,m NAC scale, where much of the scientific demand lies, are outside the runs reported here.

\paragraph{Conclusion.} Predicting one lunar instrument from another in latent space is a workable pretraining task. Its frozen features detect craters about as well as a masked-token lunar model does, gain substantially from a second input product and respond to the sun they are given. Getting there taught us more about failure than about architecture. A cross-modal JEPA can reach a very low loss by encoding where a token is instead of what the ground looks like, and loss, within-sample rank and VICReg all read this as progress. A cross-sample comparison exposes it for one cosine per target, and we would log it in any JEPA whose targets are only partly predictable from the context.
